\begin{abstract}

Vision-language models (VLMs) are increasingly deployed for video understanding, yet their ability to reason \emph{who} did \emph{what} to \emph{whom} in aggressive interactions remains largely unexamined. We introduce A$^3$Bench, a benchmark of 2,683 video clips spanning 17 aggression categories and a non-aggressive control, accompanied by over 18K multiple-choice questions organized into three tiers of compositional difficulty: basic single-attribute identification, compound two-attribute association, and detailed full aggressor-action-victim reasoning. Distractors are drawn from a structured hardness taxonomy that controls for shortcut cues including role reversal, bystander substitution, cross-video borrowing, and frequency-inverted answer construction. Zero-shot evaluation of 15 MLLMs, including both open-weight and proprietary models ranging from 7B to 72B parameters, reveals that even the best model achieves only 46.96\% accuracy, with performance degrading monotonically as questions require more compositional reasoning: 47.74\% on basic, 38.82\% on compound, and 29.58\% on detailed questions. Our analysis uncovers a systematic aggressor bias, where 37\% of true victims and 31\% of bystanders are mislabeled as aggressors, and shows that neither model scale nor explicit reasoning prompts resolve these failures. To address the identified bottleneck in role-action binding, we propose role-graph prompting, a training-free strategy that decomposes aggressive-scene understanding into a structured causal chain: detecting physical interaction, assigning interaction direction, deriving participant roles, and only then identifying the aggressive action. Role-graph prompting improves the best model's accuracy from 46.96\% to 54.32\%, with the largest gains on the compound and detailed questions that zero-shot models struggle with most.

\end{abstract}